\begin{abstract}
水产疾病诊断是一个高风险、强合规、低资源的多模态任务：养殖户通常只提供一张现场照片和一句口头主诉，而合格的诊断需要先追问关键临床字段、信息足够才下结论、无法确诊时明确拒答并建议送检，且全程不得给出具体用药剂量。本文构建 \textbf{AquaDiag}，首个面向低资源水产的多模态诊断对话基准，将诊断任务显式解耦为\textbf{病灶/体征识别}（确定性）与\textbf{病种诊断}（不确定性，含合规红线与拒答判断）两个子任务，覆盖 207 个病例、5 个病种，提供黄金追问序列、受控病灶本体（13 种病灶类型 $\times$ 10 个部位）与 6+1 项指标。在此基准上，我们发现四个代表性视觉语言模型（GPT-5.6、Qwen3-VL-8B、Qwen3.8-27B、DeepSeek-VL2）在"该追问不追问、该拒答不拒答、违规开药"上存在系统性失败：即使对话协议已显式要求先追问、图不足则送检，四者先追问率仍仅为 0\%--55.6\%，拒答准确率仅为 55.6\%--94.4\%，而合规违规率最高达 41.5\%。为修复这一问题，我们提出轻量方法 \textbf{ProDiag}，以"视觉显式化 + 可审计规则表 + 红线硬过滤"替代自由生成，在不进行任何训练的情况下，即达到诊断准确率 100\%、合规违规率 0\%、先追问率 100\%、必问字段覆盖率 94.2\% 与拒答准确率 100\%，在全部 7 项指标上超越三个基线模型。进一步的专家标注显示，53\% 的真实病灶图无法仅凭图像确定病种（爱德华氏菌病 6/6 图均需送检），印证了"病灶识别 $\neq$ 病种诊断"的天然分离，凸显可追溯证据与安全拒答在专业诊断中的必要性。
\end{abstract}
